\chapter{TỔ CHỨC HỆ QUẢN TRỊ DỮ LIỆU}\label{chap:tochuc}

Dữ liệu được chia vào bốn kho, mỗi kho phục vụ một kiểu truy vấn. SQLite là nguồn sự thật: ChromaDB, chỉ mục BM25
và đồ thị đều dựng lại được hoàn toàn từ SQLite. Cả bốn kho nằm trong một thư mục \code{data/}, nên có thể
embedding trên máy GPU rồi chép sang máy demo. Hình~\ref{fig:napdulieu} mô tả luồng nạp một PDF vào các kho.

\begin{figure}[H]
\centering
% Luồng nạp một PDF: ba lớp kiểm tra trùng trước khi ghi vào bốn kho dữ liệu.
% Lưới cột cố định (tâm cách nhau 3.8 cm), hộp rộng 3.0 cm -> không chồng nhau; tổng rộng ~14.4 cm.
\begin{tikzpicture}[
  font=\footnotesize,
  box/.style={draw, rounded corners=2pt, align=center, text width=27mm, minimum height=11mm, fill=white,
    execute at begin node=\setstretch{1.05}},
  gate/.style={box, fill=orange!12},
  store/.style={box, minimum height=9mm, fill=blue!6},
  dup/.style={box, text width=37mm, fill=black!6},
  arr/.style={-{Stealth[length=2mm]}, thick},
  lbl/.style={font=\scriptsize, fill=white, inner sep=1pt}]
  % Nhánh "đã có" (phía trên)
  \node[dup] (dup) at (7.6, 2.1) {\textbf{Đã có trong kho:} báo lý do,\\gắn bài cũ làm bài đang hỏi};
  % Hàng 1: chuỗi xử lý
  \node[box]  (pdf)   at (0.0, 0) {PDF tải lên\\(một hoặc nhiều file)};
  \node[gate] (sha)   at (3.8, 0) {\textbf{Lớp 1}\\trùng SHA-256?};
  \node[box]  (parse) at (7.6, 0) {Đọc PDF: tiêu đề,\\mã arXiv, refs};
  \node[gate] (id)    at (11.4, 0) {\textbf{Lớp 2}\\trùng mã / tiêu đề?};
  % Hàng 2: sau khi qua kiểm tra
  \node[gate] (api)   at (11.4, -2.1) {\textbf{Lớp 3}\\tra arXiv API, kiểm tra lại};
  \node[box]  (chunk) at (7.6, -2.1) {Chia đoạn\\220 từ, chồng 40};
  \node[box]  (emb)   at (3.8, -2.1) {Embedding\\+ trích khái niệm};
  % Hàng 3: bốn kho
  \node[store] (s1) at (0.0, -4.1) {SQLite};
  \node[store] (s2) at (3.8, -4.1) {ChromaDB};
  \node[store] (s3) at (7.6, -4.1) {BM25};
  \node[store] (s4) at (11.4, -4.1) {Đồ thị};
  \begin{scope}[on background layer]
    \node[draw, dashed, rounded corners=4pt, inner sep=2.5mm, fit=(s1)(s2)(s3)(s4),
      label={[font=\scriptsize\itshape]below:bốn kho dữ liệu trong thư mục \code{data/}}] (kho) {};
  \end{scope}
  % Luồng chính
  \draw[arr] (pdf) -- (sha);
  \draw[arr] (sha) -- node[lbl, above]{không} (parse);
  \draw[arr] (parse) -- (id);
  \draw[arr] (id) -- node[lbl, right]{không} (api);
  \draw[arr] (api) -- node[lbl, above]{mới} (chunk);
  \draw[arr] (chunk) -- (emb);
  \draw[arr] (emb.south) -- (emb.south |- kho.north);
  % Nhánh trùng
  \draw[arr] (sha.north) |- node[lbl, pos=0.25, left]{có} (dup.west);
  \draw[arr] (id.north) |- node[lbl, pos=0.25, right]{có} (dup.east);
\end{tikzpicture}

\caption{Luồng nạp một PDF: ba lớp kiểm tra trùng trước khi ghi vào bốn kho dữ liệu}\label{fig:napdulieu}
\end{figure}

\section{Bốn kho dữ liệu}

\begin{table}[H]
\centering\small\setstretch{1.15}
\begin{tabularx}{\linewidth}{@{}P{0.12\linewidth}P{0.24\linewidth}LL@{}}
\toprule
\textbf{Kho} & \textbf{Công nghệ} & \textbf{Nội dung} & \textbf{Phục vụ truy vấn}\\
\midrule
Metadata & SQLite (chế độ WAL), \code{scholarlens.sqlite} & Bảng \code{papers}, \code{chunks}, trạng thái pipeline & Kiểm tra trùng, lọc, chi tiết bài\\
Vector & ChromaDB 1.5 (HNSW~\cite{hnsw}, cosine), \code{chroma/} & Collection \code{chunks} (vector đoạn) và \code{papers} (vector tiêu đề + tóm tắt) & Tìm kiếm ngữ nghĩa, láng giềng tương đồng\\
Từ vựng & BM25 Okapi~\cite{bm25} (rank-bm25), trong RAM & Token của mọi đoạn đã embedding & Khớp thuật ngữ chính xác\\
Đồ thị & NetworkX~\cite{networkx} MultiDiGraph, \code{graph.json} & Nút paper, author, category, concept và 5 loại cạnh & Bài liên quan, quan hệ cho Graph RAG\\
\bottomrule
\end{tabularx}
\caption{Bốn kho dữ liệu và vai trò}\label{tab:kho}
\end{table}

\section{Lược đồ SQLite}

\begin{table}[H]
\centering\small\setstretch{1.15}
\begin{tabularx}{\linewidth}{@{}P{0.1\linewidth}LP{0.22\linewidth}@{}}
\toprule
\textbf{Bảng} & \textbf{Cột chính} & \textbf{Chỉ mục}\\
\midrule
\code{papers} & \code{paper\_id} (mã arXiv không có hậu tố phiên bản, hoặc \code{local:<sha12>}), \code{title}, \code{title\_norm}, \code{authors}, \code{abstract}, \code{categories}, \code{year}, \code{pdf\_path}, \code{file\_sha256}, \code{source}, \code{refs} (mã arXiv được trích dẫn), \code{keywords}, \code{status}, \code{n\_chunks} & \code{arxiv\_id}, \code{file\_sha256}, \code{title\_norm}, \code{status}\\
\code{chunks} & \code{chunk\_id} = \code{<paper\_id>::<idx>}, \code{kind} (summary / body), \code{section}, \code{page}, \code{text}, \code{embedded} & \code{paper\_id}, \code{embedded}\\
\bottomrule
\end{tabularx}
\caption{Lược đồ SQLite}\label{tab:sqlite}
\end{table}

Các cột dạng danh sách được lưu dưới dạng JSON. Mỗi bài đi qua các trạng thái
\code{registered} $\to$ \code{embedded} $\to$ \code{graphed}. Nhờ vậy từng module chạy lại độc lập được: module
embedding chỉ xử lý các đoạn có \code{embedded = 0}, còn bước dựng đồ thị chỉ lấy các bài đã embedding.

\section{Vector store}

Mỗi vector đoạn văn mang metadata \code{paper\_id}, \code{title}, \code{year}, \code{category}, \code{section},
\code{page}, \code{kind}. Metadata này cho phép lọc ngay trong ChromaDB, ví dụ chỉ tìm trong bài đang hỏi hoặc chỉ
lấy bài từ năm 2024 trở đi. Collection \code{papers} (một vector cho mỗi bài) dùng để tính cạnh SIMILAR\_TO: mỗi
bài nối tới tối đa 5 bài có cosine $\geq$ 0,75. Khi dựng toàn bộ đồ thị, độ tương đồng được tính chính xác bằng nhân
ma trận theo khối; khi thêm từng bài, hệ thống dùng truy vấn HNSW xấp xỉ cho nhanh.

\section{Đồ thị tri thức}

\begin{table}[H]
\centering\small\setstretch{1.15}
\begin{tabularx}{\linewidth}{@{}P{0.18\linewidth}P{0.2\linewidth}L@{}}
\toprule
\textbf{Cạnh} & \textbf{Từ $\to$ Đến} & \textbf{Cách sinh}\\
\midrule
CITES & paper $\to$ paper & Mã arXiv trong phần References của PDF; khi thêm bài mới, cạnh được nối cả hai chiều (bài mới trích dẫn bài cũ và bài cũ trích dẫn bài mới)\\
SIMILAR\_TO & paper $\to$ paper & Cosine giữa vector cấp bài, trọng số bằng cosine\\
WRITTEN\_BY & paper $\to$ author & Tác giả lấy từ arXiv API, tên chuẩn hóa làm khóa\\
MENTIONS & paper $\to$ concept & 6 khái niệm mỗi bài, trích bằng KeyBERT~\cite{keybert} rút gọn: ứng viên n-gram 1--3 từ, xếp theo cosine với vector bài, chọn đa dạng bằng MMR~\cite{mmr}\\
IN\_CATEGORY & paper $\to$ category & Chuyên mục arXiv, ví dụ cs.IR, cs.CL\\
\bottomrule
\end{tabularx}
\caption{Các loại cạnh trong đồ thị tri thức}\label{tab:canh}
\end{table}

\section{Kiểm tra trùng lặp}\label{sec:dedup}

Việc kiểm tra đi từ rẻ đến đắt và dừng ngay khi phát hiện trùng:

\begin{enumerate}[label=(\arabic*)]
\item \textbf{SHA-256 của file:} bắt lại đúng file đã tải lên trước đó, không cần đọc PDF.
\item \textbf{Mã arXiv:} đọc từ tem dọc bên lề trang 1, từ đường dẫn arxiv.org trên trang 1, hoặc từ tên file.
\item \textbf{Tiêu đề chuẩn hóa} (bỏ dấu, chữ thường, chỉ giữ chữ và số) khớp tuyệt đối.
\item \textbf{Tiêu đề gần trùng:} tỉ lệ ký tự giống nhau $\geq$ 0,92; hoặc đồng thời độ bao hàm từ
  $|A\cap B|/\min(|A|,|B|) \geq 0{,}9$ và hệ số Jaccard $|A\cap B|/|A\cup B| \geq 0{,}8$, với tiêu đề ngắn hơn có ít
  nhất 6 từ. Tiêu chí thứ hai bắt được trường hợp bỏ hoặc thêm vài từ, ví dụ ``(Extended Version)''.
\item \textbf{Kiểm tra lại sau khi tra arXiv API} với mã và tiêu đề chuẩn. Bước này bắt được PDF không có tem arXiv
  nhưng là bài đã có trong kho.
\end{enumerate}

Trên 492 bài, cơ chế này phát hiện 100\% bản trùng theo mã arXiv và 95\% tiêu đề bị biến đổi, không báo nhầm bài
thật nào (chi tiết ở Mục~\ref{sec:danhgia}).
